
\subsection{单调拉格朗日子流形}
\label{sec:2-4-4}

本小节把讨论限制到一类特殊的拉格朗日子流形。第二署名作者吴龙根在研究弗洛尔同调及其对这类拉格朗日子流形之辛拓扑的应用时，曾于 \textbf{[Oh93,96I]} 中研究这一类对象。

我们先给出 \textbf{[Oh93]} 所用的单调拉格朗日子流形定义；它正是弗洛尔在周期轨道问题中所用单调辛流形 \textbf{[Flo89I]} 的类比。令
\[
  I_\omega,I_\mu:\pi_2(M,L)\longrightarrow\mathbb R,\mathbb Z
\]
分别为第 2.4.3 小节给出的同态。

\begin{definition}{2.4.22}\label{def:2-4-22}
紧拉格朗日子流形 $L\subset(M,\omega)$ 称为\emph{单调的}\liftIndex{I053}，如果
\[
  I_\omega=\lambda I_\mu
  \qquad\text{对某个 }\lambda\geq0.
\]
我们把 $I_\mu$ 的像的正生成元定义为\emph{极小马斯洛夫数}\liftIndex{I057} $\Sigma(L)$。
\end{definition}

注意，这一情形以平凡的方式包含所谓的弱正合情形，即 $I_\omega=0$ 的情形。若忽略分次问题，它可以完全按照弗洛尔所处理的 $\pi_2(M,L)=\{0\}$ 情形来研究。因此从现在起，我们假设 $L\subset M$ 是 $\lambda>0$ 且 $I_\mu\neq0$ 的单调拉格朗日子流形。例 2.4.19 所考察的 $\mathbb RP^n\subset\mathbb CP^n$ 和克利福德环面都是单调拉格朗日子流形的例子。

正如 \textbf{[Oh93]} 所指出的，很容易看出：如果 $M$ 容许任何单调拉格朗日子流形，那么 $(M,\omega)$ 本身必为弗洛尔 \textbf{[Flo89I]} 意义下的单调辛流形。也就是说，若存在 $\lambda>0$，使得对胡雷维茨同态
$\pi_2(M)\to H_2(M;\mathbb Z)$ 像中的任意同调类 $\alpha\in H_2(M;\mathbb Z)$ 都有
\[
  c_1(\alpha)=\lambda\int_\alpha\omega,
\]
便称辛流形 $(M,\omega)$ 是单调的。

下面的命题说明，与一般拉格朗日子流形相比，这一类拉格朗日子流形具有某些特殊的辛拓扑性质。

\begin{proposition}{2.4.23}\label{prop:2-4-23}
设 $L\subset(M,\omega)$ 是单调拉格朗日子流形，则 $\Pi(L)$ 是秩为 $1$ 的自由阿贝尔群，而与之相伴的诺维科夫环 $\Lambda^R(L)$ 是一个域。
\end{proposition}

下面的定理可以用基于萨德--斯梅尔定理的维数计数来证明；证明细节请读者参阅 \textbf{[Oh93]}。

对 $J,J'\in\mathcal J_\omega$，令
\[
  \mathcal P(\mathcal J_\omega;J,J')
  =\bigl\{\{J_t\}_t\in\mathcal P(\mathcal J_\omega)
    \bigm|J_0=J,\ J_1=J'\bigr\}.
\]

\begin{theorem}{2.4.24}\label{thm:2-4-24}
\textup{（吴龙根）}\quad
设 $L_0,L_1$ 单调且极小马斯洛夫数非零。则存在子集
$\mathcal J_\omega^0(L_0,L_1)\subset\mathcal J_\omega$；并且对任意
$J_0,J_1\in\mathcal J_\omega^0(L_0,L_1)$，存在子集
\[
  \mathcal P^0(\mathcal J_\omega;J_0,J_1)
  \subset\mathcal P(\mathcal J_\omega;J_0,J_1),
\]
使得任意 $\{J_t\}_t\in\mathcal P^0(\mathcal J_\omega;J_0,J_1)$ 都满足：

\par\smallskip\noindent\textup{(2.4.25.1)}\quad
对所有 $[\ell_p,w],[\ell_q,w']\in\operatorname{Cr}(L_0,L_1)$，若
$\mu([\ell_q,w'])=\mu([\ell_p,w])$，则除非
$[\ell_p,w]=[\ell_q,w']$，模空间
$\mathcal M([\ell_p,w],[\ell_q,w'];\{J_t\}_t)$ 为空。

\par\smallskip\noindent\textup{(2.4.25.2)}\quad
对所有 $[\ell_p,w],[\ell_q,w']\in\operatorname{Cr}(L_0,L_1)$，若
$\mu([\ell_q,w'])=\mu([\ell_p,w])+1$，则模空间
$\mathcal M([\ell_p,w],[\ell_q,w'];\{J_t\}_t)$ 是紧流形，因而是有限集。

\par\smallskip\noindent\textup{(2.4.25.3)}\quad
若 $\mu([\ell_q,w'])<\mu([\ell_p,w])$，则
$\mathcal M([\ell_p,w],[\ell_q,w'];\{J_t\}_t)=\varnothing$。
\end{theorem}

该定理蕴含：对一般的 $\{J_t\}_t$，定义 2.3.12 的假设 (2.3.12.1) 与 (2.3.12.2) 成立。因此，对任意单调对子，至少可以在这样的一般 $\{J_t\}_t$ 下定义弗洛尔上边界映射 $\delta_0$。下面把覆盖变换群 $\Pi(L_0,L_1;\ell_0)$ 在诺维科夫覆盖空间 $\widetilde\Omega(L_0,L_1;\ell_0)$ 及其莫尔斯理论上的作用也计入，从而更精确地表述这一点。

另一个重要问题是模空间
$\mathcal M([\ell_p,w],[\ell_q,w'];\{J_t\}_t)$ 的相容定向。我们将在第 8 章第 8.1.3 小节证明定理 8.1.14：若对子 $(L_0,L_1)$ 相对 spin，则可以用典范方式为这一族模空间
$\mathcal M([\ell_p,w],[\ell_q,w'];\{J_t\}_t)$ 赋予相容定向。因此，如果 $(L_0,L_1)$ 还是相对 spin 对，那么在
$\mu([\ell_q,w'])=\mu([\ell_p,w])+1$ 时，整数
\[
  n([\ell_p,w],[\ell_q,w'])
  =\#\mathcal M([\ell_p,w],[\ell_q,w'];\{J_t\}_t)
\]
有良好定义。以它为矩阵元，按下式定义标准的弗洛尔上边界映射
$\delta_0:CF_R^k(L_1,L_0;\ell_0)\to CF_R^{k+1}(L_1,L_0;\ell_0)$：
\begin{equation}
  \delta_0([\ell_p,w])
  =\sum_{\mu([\ell_q,w'])=\mu([\ell_p,w])+1}
    n([\ell_p,w],[\ell_q,w'])\,[\ell_q,w'].
  \tag{2.4.26}\label{eq:2-4-26}
\end{equation}
当 $(L_0,L_1)$ 是相对 spin 对时取 $R=\mathbb Z$；一般取
$R=\mathbb Z_2$。命题 2.4.9 蕴含右端属于
$CF_R^*(L_1,L_0;\ell_0)$。

第二署名作者在 \textbf{[Oh93]} 中细化了弗洛尔 \textbf{[Flo88IV]} 的证明，证明了：对极小马斯洛夫数大于 $2$ 的单调拉格朗日子流形，在对子 $(L_0,L_1)$ 满足某个拓扑限制时有 $\delta_0\circ\delta_0=0$。一旦采用本书所定义的诺维科夫弗洛尔上链模 $CF_R^*(L_1,L_0)$，后一拓扑限制便不再需要（见定理 2.4.42）。不过，例 2.4.14 表明，即便使用诺维科夫环，也不能完全去掉对马斯洛夫指标的限制。下一小节将进一步阐明该例：把矩阵元
\[
  \bigl\langle(\delta_0\circ\delta_0)([\ell_p,w]),
    [\ell_q,w']\bigr\rangle
\]
同 $L_0$ 或 $L_1$（或二者）上的某些奇异链联系起来。当
$\mu([\ell_q,w'])-\mu([\ell_p,w])=2$ 时，紧化
$\mathcal M([\ell_p,w],[\ell_q,w'];\{J_t\}_t)$ 除断裂流线外还会出现额外边界分支；上述非零贡献正来自这些额外边界分支，如例 2.4.14 所示。

\subsection{首要障碍的出现}
\label{sec:2-4-5}
\liftIndex{I065}

暂且回到一般的对子 $(L_0,L_1)$。先描述光滑模空间
$\mathcal M^{\mathrm{reg}}([\ell_p,w],[\ell_q,w'])$ 在其紧化
$\mathcal M([\ell_p,w],[\ell_q,w'])$ 中的补集。令
\[
  \mathcal M^{\mathrm{sing}}([\ell_p,w],[\ell_q,w'])
  :=\mathcal M([\ell_p,w],[\ell_q,w'])
    \setminus\mathcal M^{\mathrm{reg}}([\ell_p,w],[\ell_q,w']).
\]
为了研究 $\mathcal M([\ell_p,w],[\ell_q,w'])$ 的边界，需要描述“余维一”层，也就是奇异轨迹
$\mathcal M^{\mathrm{sing}}([\ell_p,w],[\ell_q,w'])$ 的最高维层；把它记作
$\partial\mathcal M([\ell_p,w],[\ell_q,w'])$。在横截情形，这个层或者是一条含两个分支的断裂流线，或者是一条在边界上附着了一个圆盘的流线。

记
\begin{equation}
  \overset{\circ}{\mathcal M}_{1,0}
    ([\ell_p,w],[\ell_q,w'_0])
  \cong
  \overset{\circ}{\mathcal M}_{0,1}
    ([\ell_p,w],[\ell_q,w'_0])
  =\mathcal M^{\mathrm{reg}}([\ell_p,w],[\ell_q,w'_0]).
  \tag{2.4.27}\label{eq:2-4-27}
\end{equation}
这是定义 3.7.24 的一个特例，其中定义了
$\overset{\circ}{\mathcal M}_{k_1,k_0}
([\ell_p,w],[\ell_q,w'_0])$。分别定义映射
\[
  \operatorname{ev}:\overset{\circ}{\mathcal M}_{1,0}
    ([\ell_p,w],[\ell_q,w'_0])\longrightarrow L_1,
  \qquad
  \operatorname{ev}:\overset{\circ}{\mathcal M}_{0,1}
    ([\ell_p,w],[\ell_q,w'_0])\longrightarrow L_0
\]
如下：
\begin{equation}
  \operatorname{ev}(u)=u(0,1),
  \qquad\text{或}\qquad
  \operatorname{ev}(u)=u(0,0).
  \tag{2.4.28}\label{eq:2-4-28}
\end{equation}
这里两式依次对应上述两个映射。（当把 $u$ 看作模空间
$\overset{\circ}{\mathcal M}_{1,0}([\ell_p,w],[\ell_q,w'_0])$ 的元素时，点
$(0,1)\in\mathbb R\times[0,1]$ 被看作标记点；另一方面，当把 $u$ 看作
$\overset{\circ}{\mathcal M}_{0,1}([\ell_p,w],[\ell_q,w'_0])$ 的元素时，点
$(0,0)\in\mathbb R\times[0,1]$ 被看作标记点。）

\begin{definition}{2.4.29}\label{def:2-4-29}
设 $[\ell_p,w],[\ell_q,w']\in\operatorname{Cr}(L_0,L_1)$，且
$\beta_{(0)}\in\Pi(L_0)$、$\beta_{(1)}\in\Pi(L_1)$。对 $i=0,1$，令
\begin{align*}
&\operatorname{Int}\mathcal N
  ([\ell_p,w],[\ell_q,w']:L_1,\beta_{(1)})\\
&\quad=
  \bigcup_{\substack{\beta_{(1)},w''_1;\\
        w''_1\#\beta_{(1)}=w'}}
  \frac{
    \overset{\circ}{\mathcal M}_{1,0}
      ([\ell_p,w],[\ell_q,w''_1])_{\operatorname{ev}}
    \mathop{\times}_{\operatorname{ev}_0}
    \mathcal M_1(L_1;\beta_{(1)};J_1)}{\mathbb R},\\[4pt]
&\operatorname{Int}\mathcal N
  ([\ell_p,w],[\ell_q,w']:L_0,\beta_{(0)})\\
&\quad=
  \bigcup_{\substack{\beta_{(0)},w''_0;\\
        w''_0\#\beta_{(0)}=w'}}
  \frac{
    \overset{\circ}{\mathcal M}_{0,1}
      ([\ell_p,w],[\ell_q,w''_0])_{\operatorname{ev}}
    \mathop{\times}_{\operatorname{ev}_0}
    \mathcal M_1(L_0;\beta_{(0)};J_0)}{\mathbb R}.
\end{align*}
\end{definition}

把模空间
$\operatorname{Int}\mathcal N([\ell_p,w],[\ell_q,w']:L_1,\beta_{(1)})$
在 $\mathcal M([\ell_p,w],[\ell_q,w'])$ 中的闭包定义为
$\mathcal N([\ell_p,w],[\ell_q,w']:L_1,\beta_{(1)})$。
$\mathcal N([\ell_p,w],[\ell_q,w']:L_0,\beta_{(0)})$ 的定义类似（见图 2.4.3）。

\begin{figure}[htbp]
  \centering
  \setlength{\unitlength}{1mm}
  \begin{picture}(100,48)
    \qbezier(9,12)(47,36)(88,12)
    \qbezier(9,12)(49,-8)(88,12)
    \put(1,4){\line(1,1){8}}
    \put(1,20){\line(1,-1){8}}
    \put(88,12){\line(1,1){8}}
    \put(88,12){\line(1,-1){8}}
    \put(5,5){$p$}
    \put(88,5){$q$}
    \put(97,18){$L_0$}
    \put(97,4){$L_1$}
    \qbezier(46,45.5)(57.5,45.5)(57.5,34)
    \qbezier(57.5,34)(57.5,22.5)(46,22.5)
    \qbezier(46,22.5)(34.5,22.5)(34.5,34)
    \qbezier(34.5,34)(34.5,45.5)(46,45.5)
    \put(42,33){$\beta_{(1)}$}
    \put(46,24){\circle*{1.4}}
    \qbezier(20,10)(25,18)(30,24)
    \qbezier(31,7)(36,19)(41,28)
    \qbezier(43,5)(48,18)(52,28)
    \qbezier(56,6)(61,17)(66,24)
    \qbezier(69,9)(74,15)(78,19)
    \qbezier(38,30)(42,37)(47,45)
    \qbezier(46,23)(52,33)(56,40)
  \end{picture}
  {\renewcommand{\thefigure}{2.4.3}%
    \caption{$\mathcal N([\ell_p,w],[\ell_q,w']:L_1,\beta_{(1)})$}
    \label{fig:2-4-3}}
\end{figure}

如前所述，我们希望说
$\mathcal N([\ell_p,w],[\ell_q,w']:L_i,\beta_{(i)})$ 具有余维一。一般而言，即使对一般的 $\{J_t\}_t$，这一陈述按朴素意义也不成立（若把这些空间看成带仓西结构的空间，则一般有意义）。不过，如果可以选到这样的 $\{J_t\}_t$，并且 $(L_0,L_1)$ 单调，那么它确实成立。这里把讨论限制在单调情形。

回忆：伪全纯映射 $u:\Sigma\to M$（$\Sigma$ 是可以带边、也可以无边的黎曼曲面）称为\emph{在某处单射的}\liftIndex{I066}，如果存在 $p\in\Sigma$ 及 $p$ 的一个邻域 $U$，使得
\par\smallskip\noindent\textup{(2.4.30.1)}\quad
$\{u(p)\}\cap u(\Sigma\setminus U)=\varnothing$；

\translatornote{原书 (2.4.30.1) 把点 $u(p)$ 直接作为集合求交；译文补成单点集，以保持条件的类型。}

\par\smallskip\noindent\textup{(2.4.30.2)}\quad
$u$ 在 $p$ 处是浸入。

回忆，对一般的 $J\in\mathcal J_\omega$，有如下事实成立 \textbf{[McD87]}：

\par\smallskip\noindent\textup{(2.4.31.1)}\quad
若 $\partial\Sigma=\varnothing$，则任意伪全纯映射 $u:\Sigma\to M$ 均可分解为
$u=\overline u\circ\pi$，其中 $\pi:\Sigma\to\overline\Sigma$ 是分歧覆叠，而
$\overline u:\overline\Sigma\to M$ 在某处单射。

\par\smallskip\noindent\textup{(2.4.31.2)}\quad
若 $u:\Sigma\to M$ 在某处单射，则伪全纯映射的模空间在 $u$ 的一个邻域内光滑且具有正确维数。

（除本节外，我们不使用在某处单射性。因此，不熟悉在某处单射性的读者可以跳过本节余下部分，直接进入下一章。）

\begin{theorem}{2.4.32}\label{thm:2-4-32}
设 $L_0,L_1$ 单调，并且 $J_0,J_1\in\mathcal J_\omega^0(L_0,L_1)$、
$\{J_t\}_t\in\mathcal P^0(\mathcal J_\omega;J_0,J_1)$ 如定理 2.4.24 所述。假设
\begin{equation}
  \mu([\ell_q,w'])-\mu([\ell_p,w])\leq2.
  \tag{2.4.33}\label{eq:2-4-33}
\end{equation}
则对 $i=0,1$，
$\mathcal N([\ell_p,w],[\ell_q,w']:L_i,\beta_{(i)})$ 都是维数为
$\mu([\ell_q,w'])-\mu([\ell_p,w])-2$ 的光滑流形。
\end{theorem}

\begin{proof}
结论来自定义 2.4.29 中纤维积的简单维数计数。这里只考察 $i=1$ 的情形。

按照定义，$\mathcal N([\ell_p,w],[\ell_q,w']:L_1,\beta_{(1)})$ 是纤维积
\begin{equation}
  \frac{
    \overset{\circ}{\mathcal M}_{1,0}
      ([\ell_p,w],[\ell_q,w''_1])_{\operatorname{ev}}
    \mathop{\times}_{\operatorname{ev}_0}
    \mathcal M_1(L_1;\beta_{(1)};J_1)}{\mathbb R}
  =\frac{(\operatorname{ev}\times\operatorname{ev}_0)^{-1}(\Delta)}{\mathbb R}.
  \tag{2.4.34}\label{eq:2-4-34}
\end{equation}
这里 $\Delta\subset L_1\times L_1$ 是对角线，
$\overset{\circ}{\mathcal M}_{1,0}([\ell_p,w],[\ell_q,w''_1])$ 如 (2.4.27) 所述。回忆
\begin{align*}
  \operatorname{vir\,dim}
  \overset{\circ}{\mathcal M}_{1,0}
    ([\ell_p,w],[\ell_q,w''_1])
    &=\mu([\ell_q,w''_1])-\mu([\ell_p,w]),\\
  \operatorname{vir\,dim}\mathcal M_1(L_1;\beta_{(1)};J_1)
    &=n+\mu(\beta_{(1)})-2,
\end{align*}
\translatornote{原书此处把模空间参数写成 $\mathcal M_1(L_1,J_1;\beta_{(1)})$；译文依定义 2.4.29、紧前纤维积及本节其余用法统一次序。}
并且
\[
  \mu(\beta_{(1)})
  =\mu([\ell_q,w'])-\mu([\ell_q,w''_1]).
\]
由单调性可知，若 $\mathcal M(L_1;\beta_{(1)};J_1)\neq\varnothing$，则
$\mu(\beta_{(1)})>0$（注意 $\beta_{(1)}\neq0$）。所以 (2.4.33) 蕴含
\[
  \mu([\ell_q,w''_1])-\mu([\ell_p,w])\leq1.
\]

先考虑 $\mu([\ell_q,w''_1])-\mu([\ell_p,w])=1$ 的情形。由定理 2.4.24，模空间
\[
  \overset{\circ}{\mathcal M}_{1,0}
    ([\ell_p,w],[\ell_q,w''_1])
  \subset\mathcal M([\ell_p,w],[\ell_q,w''_1])\times\mathbb R
\]
光滑，且维数为
$\mu([\ell_q,w''_1])-\mu([\ell_p,w])=1$。

于是，$\mathcal M(L_1;\beta_{(1)};J_1)$ 的每个元素都在某处单射。事实上，如果
$v\in\mathcal M(L_1;\beta_{(1)};J_1)$ 不是在某处单射的，那么
\textbf{[KwOh00, Lazz00]} 蕴含存在 $v_i$（$i=1,\ldots,a$，$a\geq2$），使得
$\sum_i[v_i]=[v]$，且 $\mathcal M(L_1;[v_i];J_1)$ 非空。这是不可能的，因为
$\mu_{L_1}(v_i)>0$。

因此，$\mathcal M_1(L_1;\beta_{(1)};J_1)$ 光滑且维数为 $n-1$。利用在某处单射性，按如今已属标准的论证，可以证明：对一般的
$J_0,J_1,\{J_t\}_t$，评价映射
\[
  \operatorname{ev}\times\operatorname{ev}_0:
  \overset{\circ}{\mathcal M}_{1,0}
    ([\ell_p,w],[\ell_q,w''_1])
  \times\mathcal M_1(L_1;\beta_{(1)};J_1)
  \longrightarrow L_1\times L_1
\]
横截于 $\Delta$。所以模空间 (2.4.34) 成为光滑流形，其维数为
\begin{align}
 &\mu([\ell_q,w''_1])-\mu([\ell_p,w])+n+\mu(\beta_{(1)})-2-n\notag\\
 &\quad=\mu([\ell_q,w''_1])-\mu([\ell_p,w])+\mu(\beta_{(1)})-2\notag\\
 &\quad=\mu([\ell_q,w'])-\mu([\ell_p,w])-2.
 \tag{2.4.35}\label{eq:2-4-35}
\end{align}
这里末行按定义 2.4.29 的关系 $w''_1\#\beta_{(1)}=w'$ 写出。\translatornote{原书 (2.4.35) 末行印作 $w'_1$；同页定理的维数陈述以及定义 2.4.29 的粘接关系均表明此处应为目标提升 $w'$。}

接着考虑 $\mu([\ell_q,w''_1])-\mu([\ell_p,w])\leq0$。假设 (2.4.34) 非空。于是定理 2.4.24 蕴含
$[\ell_q,w''_1]=[\ell_p,w]$，而且
$\overset{\circ}{\mathcal M}_{1,0}([\ell_p,w],[\ell_q,w''_1])$ 必须由定常元素
$u\equiv p$ 构成。再者，若 $(u,v)$ 是 (2.4.34) 的元素，其中
$u\in\overset{\circ}{\mathcal M}_{1,0}([\ell_p,w],[\ell_q,w''_1])$ 且
$v\in\mathcal M_1(L_1;\beta_{(1)};J_1)$，则 $v$ 必须通过 $p$。此外，
$\mu(\beta_{(1)})\leq2$。

若 $v$ 不是在某处单射的，则 $v$ 的像包含在若干映射
$v_i:(D^2,\partial D^2)\to(M,L_1)$ 的像之并中，而且
$\mu([v_i])<\mu([v])\leq2$（\textbf{[KwOh00, Lazz00]}）。因此
$\mathcal M_1(L_1;[v_i];J_1)$ 是维数至多为 $n-1$ 的紧光滑流形。故在必要时重新选择 $J_1$ 后，评价映射
$\operatorname{ev}:\mathcal M_1(L_1;[v_i];J_1)\to L_1$ 的像可以避开零维集合
$L_0\cap L_1$。

所以，只要 $(u,v)$ 是 (2.4.34) 的元素，就可以假设 $v$ 总是在某处单射的。于是
$\mathcal M_1(L_1;\beta_{(1)};J_1)$ 是光滑流形，纤维积 (2.4.34) 横截。定理 2.4.32 得证。
\end{proof}

由此得到：

\begin{theorem}{2.4.36}\label{thm:2-4-36}
设 $L_0,L_1$ 和 $\{J_t\}_t$ 如定理 2.4.32 所述，并假设 (2.4.33) 成立，则有恒等式
\begin{align*}
  \partial\mathcal M([\ell_p,w],[\ell_q,w'])
  ={}&\bigcup_{\mu([\ell_r,w''])=\mu([\ell_q,w'])-1}
    \\[-2pt]
  &\quad
    \mathcal M([\ell_p,w],[\ell_r,w''])
    \times\mathcal M([\ell_r,w''],[\ell_q,w'])\\
  &{}+\bigcup_{\beta_{(0)}}
    \mathcal N([\ell_p,w],[\ell_q,w']:L_0,\beta_{(0)})\\
  &{}+\bigcup_{\beta_{(1)}}
    \mathcal N([\ell_p,w],[\ell_q,w']:L_1,\beta_{(1)}).
\end{align*}
\end{theorem}

\begin{proof}
令 $u_i$ 是 $\mathcal M([\ell_p,w],[\ell_q,w'])$ 中的一个发散序列。由格罗莫夫紧性定理，在必要时取子序列后，极限中必发生下列情形之一：

\par\smallskip\noindent\textup{(2.4.37.1)}\quad
$u_i$ 分裂成若干连接轨道之和。

\par\smallskip\noindent\textup{(2.4.37.2)}\quad
在 $\mathbb R\times[0,1]$ 的某些点发生球面冒泡。

\par\smallskip\noindent\textup{(2.4.37.3)}\quad
在 $\mathbb R\times[0,1]$ 的某些边界点，即 $\mathbb R\times\{0,1\}$ 中的某些点，发生圆盘冒泡。

(2.4.37.1) 对应定理 2.4.36 的第一项，(2.4.37.3) 对应其中的第二、第三项。(2.4.37.2) 是余维 $2$ 的现象，因而不会出现在余维 $1$ 的边界中。最后这个维数计数由定理 2.4.32 保证，故得证。
\end{proof}

由标准的配边论证，定理 2.4.36 蕴含恒等式
\begin{align}
  0={}&\bigl\langle(\delta_0\circ\delta_0)([\ell_p,w]),[\ell_q,w']\bigr\rangle
  +\sum_{\beta_{(0)}}
    \#\mathcal N([\ell_p,w],[\ell_q,w']:L_0,\beta_{(0)})\notag\\
  &{}+\sum_{\beta_{(1)}}
    \#\mathcal N([\ell_p,w],[\ell_q,w']:L_1,\beta_{(1)}).
  \tag{2.4.38}\label{eq:2-4-38}
\end{align}
这里
$\langle(\delta_0\circ\delta_0)([\ell_p,w]),[\ell_q,w']\rangle$
表示 $[\ell_q,w']$ 在 $(\delta_0\circ\delta_0)([\ell_p,w])$ 中的系数；
$\#$ 表示带符号计数。这里不解释符号（即模空间的定向），因为后续各章、特别是第 8 章会详细讨论它。

现在像 \textbf{[Oh93,96I]} 那样，分别研究如下两种情形：

\par\smallskip\noindent\textup{(2.4.39.1)}\quad
$L_0,L_1$ 的极小马斯洛夫数大于 $2$；

\par\smallskip\noindent\textup{(2.4.39.2)}\quad
$L_0,L_1$ 之一或二者的极小马斯洛夫数等于 $2$。

先考虑情形 (2.4.39.1)。如果 $\beta_{(i)}\in\Pi(M;L_i)$ 满足
$\mu(\beta_{(i)})>2$，那么当
$\mu([\ell_q,w'])-\mu([\ell_p,w])=2$ 时，对一般的 $\{J_t\}_t$ 有
$\mathcal N([\ell_p,w],[\ell_q,w']:L_i,\beta_{(i)})=\varnothing$。基于这一观察，第二署名作者在 \textbf{[Oh93]} 中用维数计数证明了 (2.4.38) 的额外项均为零。下面用本书所采用的较新语言回忆该论证。

只考虑 $i=1$。\textbf{[Oh93]} 考察了模空间
\[
  \mathcal M^\partial(L_1;\beta_{(1)})
  =\bigl(\widetilde{\mathcal M}(L_1;\beta_{(1)})
    \times S^1\bigr)/\operatorname{PSL}(2;\mathbb R)
\]
和评价映射
\[
  \operatorname{ev}^\partial:
  \mathcal M^\partial(L_1;\beta_{(1)})\longrightarrow L_1,
  \qquad
  \operatorname{ev}^\partial([w,t])=w(t).
\]
注意，$\mathcal M^\partial(L_1;\beta_{(1)})$ 的维数是
$\mu(\beta_{(1)})+n-2$。还要注意，
$\mathcal M^\partial(L_1;\beta_{(1)})$ 实际上等于
$\mathcal M_1(L_1;\beta_{(1)})$，而 $\operatorname{ev}^\partial$ 等于
$\operatorname{ev}_0:\mathcal M_1(L_1;\beta_{(1)})\to L_1$。因此，如果假设
$L_0,L_1$ 的极小马斯洛夫数大于 $2$，从而 $\mu(\beta_{(1)})>2$，那么在
$\mathcal M_1(L_1;\beta_{(1)})$ 横截且维数正确时，有
\[
  \dim\mathcal M_1(L_1;\beta_{(1)})>\dim L_1.
\]
在这种情况下，链
$\operatorname{ev}_0:\mathcal M_1(L_1;\beta_{(1)})\to L_1$ 是退化的，因而作为
$L_1$ 上维数为 $\dim\mathcal M_1(L_1;\beta_{(1)})$ 的链等于零；这正是下述论证背后的原因。

回忆，只有当 $\mu([\ell_q,w'])-\mu([\ell_p,w])=2$ 时，
$\langle(\delta_0\circ\delta_0)([\ell_p,w]),[\ell_q,w']\rangle$ 才可能非零。因此，由纤维积
\begin{equation}
  \frac{
    \overset{\circ}{\mathcal M}_{1,0}
      ([\ell_p,w],[\ell_q,w''_1])_{\operatorname{ev}}
    \mathop{\times}_{\operatorname{ev}_0}
    \mathcal M_1(L_1;\beta_{(1)})}{\mathbb R}
  \tag{2.4.40}\label{eq:2-4-40}
\end{equation}
的维数公式 (2.4.35) 以及定理 2.4.24，可知其维数为
$\mu([\ell_q,w''_1])-\mu([\ell_p,w])+\mu(\beta_{(1)})-2=0$。所以，由假设 (2.4.39.1)，有
\begin{equation}
  \mu([\ell_q,w''_1])-\mu([\ell_p,w])
  =2-\mu(\beta_{(1)})<0.
  \tag{2.4.41}\label{eq:2-4-41}
\end{equation}
(2.4.41) 与 (2.4.25.3) 蕴含
\[
  \overset{\circ}{\mathcal M}_{1,0}
    ([\ell_p,w],[\ell_q,w''_1])=\varnothing,
\]
故 $\mathcal N([\ell_p,w],[\ell_q,w']:L_1,\beta_{(1)})$ 必为空。这证明了第二署名作者在 \textbf{[Oh93]} 中得到的如下定理（下面的陈述略有加强；例如，当 $(L_1,L_0)$ 是相对 spin 对时，可以取 $R=\mathbb Z$）。

\begin{theorem}{2.4.42}\label{thm:2-4-42}
\textup{（吴龙根）}\quad
设 $L_0,L_1$ 是单调拉格朗日子流形，并假设它们的极小马斯洛夫数大于 $2$。则由 (2.4.26) 定义的上边界算子
\[
  \delta_0:CF_R^*(L_1,L_0;\ell_0)
  \longrightarrow CF_R^{*+1}(L_1,L_0;\ell_0)
\]
满足
\[
  \delta_0\circ\delta_0=0.
\]
这里，当 $(L_0,L_1)$ 是相对 spin 对时取 $R=\mathbb Z$；一般取
$R=\mathbb Z_2$。
\end{theorem}

因此，在情形 (2.4.39.1) 中，定义弗洛尔上同调
\[
  HF(L_1,L_0;\ell_0)=\operatorname{Ker}\delta_0/
  \operatorname{Im}\delta_0.
\]
可以证明，它只依赖于 $L_0,L_1$ 的哈密顿同痕类。这里不写出细节，因为第 5 章将证明更一般的结果。

现在研究情形 (2.4.39.2)，即 $L_1$ 或 $L_0$ 的极小马斯洛夫数为 $2$。此时 (2.4.41) 变为
\[
  \mu([\ell_q,w''_1])-\mu([\ell_p,w])
  =2-\mu(\beta_{(1)})\leq0,
\]
于是必有
$\mu([\ell_q,w''_1])-\mu([\ell_p,w])=0$ 且
$\mu(\beta_{(1)})=2$。由定理 2.4.24，前一个等式蕴含 $p=q$，并且若
$(u,v)$ 是 (2.4.40) 的元素，则 $v$ 必须通过 $p$。再者，链
$\operatorname{ev}_0:\mathcal M_1(L_1;\beta_{(1)})\to L_1$ 定义一个 $n$ 维紧闭链，因而是基本类 $[L_1]$ 的某个倍数。如例 2.4.14 所示，这个闭链确实可能非零，故相应的纤维积也确实可能存在。结果是：即使对那些只需扰动殆复结构便可获得大部分所需横截性的单调拉格朗日子流形，$\delta_0\circ\delta_0\neq0$ 也确实可能发生。

如果对某个哈密顿微分同胚 $\psi$ 有 $L_1=\psi(L_0)$，那么在较弱的假设——极小马斯洛夫数不小于 $2$——下就有 $\delta_0\circ\delta_0=0$。见 \textbf{[Oh93]} 或本书第 3.6.3 小节。

为了在定理 2.4.42 的情形之外定义拉格朗日子流形的弗洛尔上同调，本书引入了一个重要想法：对 $L_i$ 施加如下同调条件。当 $\beta_{(i)}$ 具有最小可能的辛面积时，链
$\operatorname{ev}_0:\mathcal M_1(L_i;\beta_{(i)})\to L_i$ 成为闭链；要求它同调于零，然后用以它为边界的一条链把标准弗洛尔上边界映射 $\delta_0$ 修正为 $\delta$，从而使 $\delta\circ\delta=0$。事实证明，要相容地解释这种参数依赖并同时处理横截性问题，需要新的同调代数和滤过 $A_\infty$ 代数形变理论中的大量代数工具。

\begin{remark}{2.4.43}\label{rem:2-4-43}
使用第 3.7 节引入的滤过 $A_\infty$ 双模术语，定理 2.4.36 可以写成
\begin{align}
  \bigl\langle(\mathfrak n_{0,0}\circ\mathfrak n_{0,0})
    ([\ell_p,w]),[\ell_q,w']\bigr\rangle
  ={}&\mathbin{\pm}\bigl\langle
    \mathfrak n_{1,0}(\mathfrak m_0(1),[\ell_p,w]),[\ell_q,w']
    \bigr\rangle\notag\\
  &{}\mathbin{\pm}\bigl\langle
    \mathfrak n_{0,1}([\ell_p,w],\mathfrak m_0(1)),[\ell_q,w']
    \bigr\rangle.
  \tag{2.4.44}\label{eq:2-4-44}
\end{align}
也就是说，(2.4.38) 的第二、第三项是 (2.4.44) 的右端，而 (2.4.38) 的第一项是 (2.4.44) 的左端。注意，(2.4.44) 来自恒等式
\[
  (\mathfrak n_{0,0}\circ\mathfrak n_{0,0})(x)
  \mathbin{\pm}\mathfrak n_{1,0}(\mathfrak m_0(1),x)
  \mathbin{\pm}\mathfrak n_{0,1}(x,\mathfrak m_0(1))=0,
\]
这是滤过 $A_\infty$ 双模定义方程的一个特例。确切地说，它是
\[
  \operatorname{Hom}(D,D)
  =\operatorname{Hom}(B_0C_1[1]\otimes D\otimes B_0C_0[1],D)
\]
这一分量上的 $\widehat d^2=0$。见定义 3.7.5 (1)。
\end{remark}

\begin{remark}{2.4.45}\label{rem:2-4-45}
最后说明带基点道路 $\ell_0$ 在本书构造中的作用。我们的拉格朗日交理论应当建立在带基点的道路空间上，而非无基点的道路空间上；这一观点在定义拉格朗日子流形哈密顿同痕下的链同态时很重要（见第 5.3 节）。为了给弗洛尔复形赋予绝对分次，并在哈密顿同痕下定义诺维科夫环上弗洛尔上同调之间的保次数链同态，还需要选取 $\ell_0^*\Lambda(M)$ 的一个截面 $\lambda^0$（见第 2.2.2 小节）。关于这一点的进一步说明见第 5.1 节。
\end{remark}
